\chapter*{TÓM TẮT}
\addcontentsline{toc}{chapter}{TÓM TẮT}
Tính cấp thiết của nghiên cứu xuất phát từ số lượng ngày càng tăng các bệnh của hệ mạch máu não, đặc biệt là đột quỵ, một trong những nguyên nhân hàng đầu gây tàn tật và tử vong trên thế giới. Để cải thiện các mô hình dự đoán nguy cơ đột quỵ về cả hiệu quả và khả năng diễn giải, nhóm tác giả đề xuất tích hợp các thuật toán học máy hiện đại với các phương pháp giảm chiều dữ liệu, cụ thể là XGBoost và phân tích thành phần chính (PCA) được tối ưu hóa, qua đó cấu trúc dữ liệu và tăng tốc độ xử lý, đặc biệt đối với các tập dữ liệu lớn. Trí tuệ nhân tạo giải thích được (XAI), cụ thể là SHAP, được tích hợp vào quá trình PCA, làm tăng tính minh bạch và khả năng diễn giải, giúp các chuyên gia y tế hiểu rõ hơn các yếu tố nguy cơ.

\noindent Mô hình được kiểm thử trên hai tập dữ liệu với pipeline chống rò rỉ dữ liệu. ROC-AUC kiểm định chéo đạt 0.824 (Dataset 1) và 0.997 (Dataset 2). Trên tập kiểm tra giữ nguyên tỷ lệ mất cân bằng gốc: Dataset 1 đạt accuracy 78.2\% (recall 78\%, precision 15.5\%, MCC 0.282); Dataset 2 đạt accuracy 96.6\% (recall 99.4\%, precision 55.9\%, MCC 0.732). PCA cải thiện có ý nghĩa thống kê ở Dataset 2 nhưng làm giảm nhẹ độ chính xác ở Dataset 1.

\noindent Kết quả cho thấy việc kết hợp XGBoost, PCA song song hóa bằng OpenMP và SHAP là khả thi và khả năng diễn giải được cho chuyên gia y tế. Đồng thời, kết quả cũng cho thấy rõ giới hạn thực tế của bài toán dự đoán đột quỵ trên dữ liệu mất cân bằng mạnh: độ chính xác tổng thể cao có thể gây hiểu lầm nếu không xem xét cùng với precision, recall, ROC-AUC và MCC. Do đó, phương pháp đề xuất có tiềm năng ứng dụng trong sàng lọc và hỗ trợ ra quyết định lâm sàng, với điều kiện được diễn giải đúng cách và kết hợp với đánh giá chuyên môn.

\noindent\textbf{Từ khóa:} Học máy, công nghệ tính toán song song, phương pháp SHAP, cân bằng lớp, phương pháp PCA.
